\section{Future Directions}

Several extensions are natural candidates for future work:

\textbf{Dynamic occupancy models.} The static occupancy framework developed
here can be extended to a dynamic formulation in which colonisation and
extinction probabilities are modelled as functions of environmental change,
enabling formal inference about range shifts under climate scenarios.

\textbf{Integration with process-based vegetation models.} Linking the
statistical models to mechanistic simulators such as LPJ-GUESS would enable
data assimilation and improved projections of forest dynamics under alternative
management and climate trajectories.

\textbf{Approximate inference for scalability.} Variational Bayes and INLA
approaches warrant systematic benchmarking against the gold-standard HMC
results, particularly for routine monitoring applications that require
frequent model updates.

\textbf{Decision-theoretic design.} The Bayesian framework naturally
accommodates optimal experimental design: future monitoring effort can be
allocated to sites and species that maximise expected information gain,
formalising the adaptive feedback loop depicted in
Figure~\ref{fig:framework}.

The methods developed in this thesis are intended to be practical tools for the
ecologists, foresters, and conservation planners who work daily to understand
and protect Australia's unique forest heritage. All code, data-processing
scripts, and supplementary materials are archived in a public repository to
facilitate replication and extension by the research community.

% ===========================  BACKMATTER  ==================================
\backmatter

\begin{thebibliography}{99}

\bibitem{gelman2013}
A.\ Gelman, J.\ B.\ Carlin, H.\ S.\ Stern, D.\ B.\ Dunson, A.\ Vehtari, and
D.\ B.\ Rubin,
\textit{Bayesian Data Analysis}, 3rd ed.\ Boca Raton, FL: CRC Press, 2013.

\bibitem{elith2009}
J.\ Elith and J.\ R.\ Leathwick,
\enquote{Species distribution models: Ecological explanation and prediction
across space and time,}
\textit{Annu.\ Rev.\ Ecol.\ Evol.\ Syst.}, vol.\ 40, pp.\ 677--697, 2009.

\bibitem{keith2014}
H.\ Keith, D.\ B.\ Lindenmayer, B.\ G.\ Mackey, D.\ Blair, L.\ Carter,
L.\ McBurney, S.\ Okada, and T.\ Konishi-Nagano,
\enquote{Managing temperate forests for carbon storage: Impacts of logging
versus protection on carbon stocks,}
\textit{Ecosphere}, vol.\ 5, no.\ 6, art.\ 75, pp.\ 1--29, 2014.

\bibitem{bradshaw2012}
C.\ J.\ A.\ Bradshaw,
\enquote{Little left to lose: Deforestation and forest degradation in
Australia since European colonization,}
\textit{J.\ Plant Ecol.}, vol.\ 5, no.\ 1, pp.\ 109--120, 2012.

\bibitem{hoffman2014}
M.\ D.\ Hoffman and A.\ Gelman,
\enquote{The No-U-Turn Sampler: Adaptively setting path lengths in Hamiltonian
Monte Carlo,}
\textit{J.\ Mach.\ Learn.\ Res.}, vol.\ 15, no.\ 1, pp.\ 1593--1623, 2014.

\end{thebibliography}
